\section{Model and numerical protocol}

The common model is a spinless nearest-neighbor tight-binding Hamiltonian on a
square lattice,
\begin{equation}
    H = \sum_i \epsilon_i c_i^\dagger c_i
        + \sum_{\langle i,j\rangle} t
        \left(c_i^\dagger c_j+c_j^\dagger c_i\right),
\end{equation}
with onsite energy $\epsilon_i=0$ and hopping $t=-1$. The bottom of the
infinite square-lattice band is at $E=-4$, so low-energy spectral quantities are
reported with the kinetic convention
\begin{equation}
    E_{\mathrm{kin},k}=E_k+4 .
\end{equation}
Closed-dot spectra are computed by direct sparse diagonalization of finalized
Kwant Hamiltonians~\cite{groth2014kwant,lehoucq1998arpack,virtanen2020scipy}.
The same tight-binding convention underlies the transport diagnostic, where
two-terminal spinless conductance is written as $G=T$ in Landauer units, with
lead self-energies handled by Kwant~\cite{datta2005quantumtransport}.
For two conductance curves sampled on the same energy grid, the raw curve
distance used in the transport diagnostic is the normalized RMS distance
\begin{align}
    R_F[X] &= \left[N_f^{-1}\sum_{i\in F}X(E_i)^2\right]^{1/2}, \\
    D_{\mathrm{raw}}(G_a,G_b)
    &= \frac{R_F[G_a-G_b]}
    {\left(R_F[G_a]^2+R_F[G_b]^2\right)^{1/2}},
    \label{eq:transport-distance}
\end{align}
where $F$ is the set of energy samples where both curves are finite and
$N_f=|F|$. The contact-width effect is
$D_{\mathrm{raw}}(G_{n=2,W=4},G_{n=2,W=10})$. The shape effect at fixed
$W$ is the mean pairwise $D_{\mathrm{raw}}$ among the tested
$n=\{1.2,2.0,4.0\}$ curves at that $W$. A separate shift/rescale diagnostic
minimizes the same distance after integer index shifts up to $12$ samples and
an affine rescaling of one curve; those values are reported only as a diagnostic
check. The plotted transport bars use the raw normalized distances in
Eq.~(\ref{eq:transport-distance}).

The geometry family is the lattice realization of a superellipse,
\begin{equation}
    \left|\frac{x}{a}\right|^n+
    \left|\frac{y}{b}\right|^n \leq 1,
    \qquad b=a r_{\mathrm{AR}},
\end{equation}
where $a$ sets the major size, $r_{\mathrm{AR}}$ is the aspect ratio, and $n$
controls the boundary shape. The completed closed-spectrum benchmark contains
$140$ geometries with
$n=\{1.2,2.0,3.0,4.0\}$,
$a=\{24,27,30,33,36\}$, and
$r_{\mathrm{AR}}\in\{0.67,0.72,0.78,0.83,0.89,0.94,1.0\}$.
The primary spectral targets in the benchmark analysis are $E_0$ and
$dE_1=E_1-E_0$; higher gaps are retained as diagnostics because degeneracies
and level ordering make them less stable targets.

The analysis uses only pre-existing outputs listed in the source inventory. All
numerical values used in the audit are traceable to versioned repository outputs
there; the inventory separates confirmed numerical values and file provenance
from the interpretation used in this manuscript. The manuscript does not
introduce additional parameter scans beyond these outputs. No new physics
calculations were performed during manuscript drafting. The evidence is
grouped by protocol:

\begin{enumerate}
    \item a closed-spectrum benchmark dataset and Ridge versus small-MLP
    comparison;
    \item one-shot Q inverse screening and preregistered S-objective screening;
    \item a direct-Kwant magnetic-response ranking check;
    \item a strict continuum-residual protocol, including a circular
    Bessel anchor and finite-difference reference checks;
    \item a tight-binding-only self-scaling diagnostic after the FD reference became
    insufficient for shape-sensitive residuals;
    \item an open-transport contact-width diagnostic.
\end{enumerate}

\paragraph*{Data and code availability.}
The numerical values summarized in this paper are taken from versioned outputs
in the public project repository at
\url{https://github.com/Pururin-ux/Diplom-Physcis-ML}. The source inventory is a
provenance table that maps each manuscript value to the corresponding branch or
commit, report path, and CSV or report source. A repository release and archived
DOI should be provided for submission. DOI to be added after archival release.

The shared analysis rule is baseline-first. Before interpreting an observable
as shape physics or inverse-design evidence, the observable must be compared
against the simplest competing explanations: confinement scale, area, perimeter,
aspect ratio, isotropy, symmetry and degeneracy, lattice boundary realization,
reference-solver uncertainty, and, for transport, lead-dot coupling. Machine
learning is treated as a candidate generator or surrogate approximation tool,
not as physical truth.
